\documentclass[11pt]{article}
\usepackage{amsmath}
\title{A shared projection error does not create independent monitor hits}
\date{26 November 2026}
\begin{document}
\maketitle

\paragraph{Source boundary.}
At Navigation2 commit \texttt{346bd05f40b62ae60b96ba86e02eb491c8dba444},
the selected polygon counts returns inside its boundary, compares
the count with \texttt{min\_points}, and updates consecutive trigger
and release counters. The replay compiles these selected unchanged
methods and the original point-in-polygon helper against a serial
fixture. It uses one fixed rectangle and a single logical list of
projected points, not the ROS node, arbitration, motion, or timing.
For the unit rectangle in the fixture the source accepts
$0\leq x<1$ and $0<y\leq1$; boundary treatment matters.

\paragraph{Law and scope.}
Suppose the uncertainty in all the listed projected returns in an
observation epoch is the same additive $x$ shift $z$ in a warranted
closed interval. The set of ordinates and the rectangle are fixed
within each evaluation; point lists can vary by evaluation. This is
a scalar model for one shared uncertainty in the input projection,
motivated by the two-time source in the preceding study, not a claim
that its native TF implementation emits epochs or guarantees this law.
Between epochs one may supply a new bound $[a,b]$ on the absolute
shift and an interval $[d_-,d_+]$ on its increment. Every combination
compatible with these constraints is admitted. Latch state and
counters persist across the boundary. The conditional exact result
does not handle ambiguous event ordering, rotations or correlated
changes that are not described by these declared transitions.

\paragraph{Exact history interface.}
Represent the relation between $z$ and the possible latch/counter
states as a piecewise-constant map, with separate cells for boundary
points and the open intervals between them. Initially put the
admitted initial states at all allowed shifts. Each return $(x,y)$
with admitted $y$ contributes on $[-x,1-x)$, clipped to the support.
Split at all the new endpoints, and apply the deterministic
hysteresis update to every prior state at that cell's hit count.
Old boundaries are discarded only if the point and both adjacent
open cells have identical state sets. Counters can be saturated at
their respective thresholds without changing a future observable
latch transition: stepping above a threshold is indistinguishable
until the opposite hit resets it. This excludes machine integer
overflow; the native comparison uses bounded traces.

At a rebase, a prior singleton $u$ can contribute its states on the
closed interval $[u+d_-,u+d_+]$, and an open cell $(u,v)$ can contribute
on $(u+d_-,v+d_+)$. These are exact sumsets with the increment interval.
Intersect their union of state contributions with $[a,b]$, retaining
each contributing state on every target cell. A justified new
independent interval admits the Cartesian product of the former
state set and new shift interval. It does not restart the counters.
The forward relation is exact by induction: splitting preserves every
prior history, deterministic update advances its state, and each
target-state contribution at an epoch boundary has a legal predecessor
shift. A missing contribution never provides affirmative evidence;
an empty relation is inconsistent with the assumed envelopes.

The incremental computation requires no later scans. A complete
direct account using the same joint error law can recover the
same decisions. The independent replay compares full shift/state
relations against offline enumeration within an epoch and intersects
preimages for bounded two-cut rebases. It separately exercises the
native hit/counter logic on 2{,}290 sampled frames. The symbolic
algorithm uses exact rational boundary locations, while the C++
test uses exactly representable dyadic coordinates.

\paragraph{Sharp two-cut witnesses.}
Start inactive with one point sufficient to hit and two consecutive
hits required to trigger. Under shifts in $[-1,1]$, the respective
returns $x=-3/4$ and $x=7/4$ can each hit, but their hit sets do not
intersect. The latch never turns on under one common shift. At a
bounded rebase they can both hit exactly when the symmetric permitted
increment is strictly larger than $3/2$: the first requires
$z_0\ge3/4$, the second $z_1<-3/4$.
Conversely, start active with two consecutive misses required to
release. For returns $x=1/4$ and $x=3/4$ and shifts in
$[-1/2,1/2]$, each can miss, but release is impossible under one
common shift. The misses require $z_0<-1/4$ and $z_1\ge1/4$,
so release only becomes possible at increments strictly above $1/2$.
Both counterexamples keep $y$ strictly inside. The replay checks
the strict boundary. A per-frame maximum is a valid upper envelope
for its hit count but does not tell the monitor whether those hits
can occur on one history; the opposite error arises for misses.

The result is about preserving the configured monitor predicate, not
certifying that this predicate ensures obstacle detection or clearance.
New range or motion-reference evidence costs work and its validity at
the actuation cut still needs a qualified source and explicit authority.
No plant, hardware, new architecture disposition, or mission acceptance
is claimed.
\end{document}
